\capitulo{7}{Conclusiones y Líneas de trabajo futuras}

El desarrollo de este Trabajo Fin de Grado ha permitido \textbf{diseñar, implementar y desplegar un sistema completo} basado en la técnica \eng{Retrieval-Augmented Generation} (RAG) aplicado a la consulta inteligente de licitaciones públicas. El proyecto ha demostrado que la combinación de modelos de lenguaje de gran tamaño con sistemas de recuperación semántica permite superar gran parte de las limitaciones tradicionales de los LLM, especialmente aquellas relacionadas con el conocimiento estático, la falta de acceso a información específica y las alucinaciones generadas por el modelo.

Uno de los principales logros del proyecto ha sido la construcción de un \eng{pipeline} completo de indexación documental capaz de automatizar la obtención de licitaciones públicas mediante técnicas de \eng{Web Scraping}, procesar los documentos PDF, fragmentarlos en \eng{chunks}, generar \eng{embeddings} y almacenarlos en una base de datos vectorial Qdrant. Gracias a ello, el sistema es capaz de \textbf{recuperar información contextual relevante y utilizarla para enriquecer el contexto} del modelo de lenguaje.

A nivel técnico, el proyecto ha permitido profundizar en múltiples áreas de la Ingeniería Informática relacionadas con la inteligencia artificial, el procesamiento del lenguaje natural, las bases de datos vectoriales, el desarrollo \eng{web} y el despliegue de aplicaciones en producción. La implementación de la aplicación \eng{web} mediante Flask, junto con el uso de Docker, Gunicorn y Nginx, ha permitido construir un entorno reproducible, escalable y preparado para un despliegue real. Además, la integración de autenticación de usuarios, gestión documental, internacionalización, soporte \eng{responsive} y paneles de administración ha permitido transformar el prototipo inicial en una \textbf{aplicación funcional} y utilizable.

Otro aspecto especialmente relevante ha sido la incorporación de mecanismos automáticos de evaluación del sistema RAG mediante \eng{frameworks} como RAGAS y ARES. Esto ha permitido analizar métricas relacionadas con la relevancia, fidelidad y precisión de las respuestas generadas, aportando una base objetiva para comparar configuraciones y mejorar el comportamiento del sistema. La evaluación ha puesto de manifiesto que la calidad final de las respuestas no depende únicamente del modelo LLM utilizado, sino también del proceso de \eng{chunking}, de los \eng{embeddings}, de la calidad de los documentos y de la estrategia de recuperación semántica empleada.

Durante el desarrollo del proyecto también se han identificado diversas dificultades técnicas. Una de las más importantes ha sido la gran heterogeneidad de formatos presentes en los documentos PDF de licitaciones, especialmente en el proceso de conversión automática a Markdown y detección de secciones. Esto ha demostrado que, aunque existen herramientas avanzadas de OCR y procesamiento documental, la extracción estructurada de información sigue siendo uno de los retos más complejos en sistemas RAG aplicados a documentación real.

Por otro lado, el proyecto ha permitido comprobar la importancia de una buena planificación incremental basada en metodologías ágiles como \eng{Scrum}. La organización del desarrollo en \eng{sprints} facilitó la adaptación continua del proyecto ante problemas técnicos, cambios en la página sobre la que se realiza el \eng{Web Scraping} o modificaciones en las decisiones arquitectónicas tomadas inicialmente.

Desde un punto de vista personal y académico, este trabajo ha supuesto una oportunidad para aplicar de forma práctica numerosos conocimientos adquiridos durante el grado, así como para profundizar en tecnologías emergentes relacionadas con los LLM y los sistemas RAG. Además, el proyecto ha permitido adquirir experiencia en investigación, diseño de arquitecturas \eng{software}, despliegue de aplicaciones, evaluación de sistemas de inteligencia artificial y documentación técnica de proyectos complejos.

\section{Líneas de trabajo futuras}

Aunque el sistema desarrollado cumple los objetivos planteados inicialmente, existen múltiples líneas de mejora y ampliación que podrían desarrollarse en trabajos futuros.

Una de las mejoras más importantes sería \textbf{optimizar el \eng{pipeline} de recuperación semántica} incorporando técnicas avanzadas de \eng{chunking} semántico y modelos de \eng{reranking}. Actualmente, el sistema emplea recuperación semántica basada en similitud coseno sobre los \eng{embeddings} almacenados en Qdrant, junto con filtros híbridos por número de expediente y tipo documental. Además, los fragmentos recuperados se ordenan según su puntuación de similitud antes de incorporarse al \eng{prompt} del LLM. Sin embargo, la arquitectura no incorpora todavía una etapa de \eng{reranking} semántico especializada. Como línea de trabajo futura, podría integrarse un modelo \eng{Cross-Encoder} capaz de reevaluar y reordenar los \eng{chunks} recuperados teniendo en cuenta conjuntamente la consulta del usuario y el contenido completo de cada fragmento. Esto permitiría mejorar significativamente la precisión contextual y la relevancia de las respuestas generadas, especialmente en consultas complejas, ambiguas o sobre documentos extensos donde múltiples fragmentos presentan similitudes vectoriales próximas.

Otra línea de desarrollo futuro podría ser incorporar \textbf{memoria conversacional} persistente. Actualmente, las consultas se procesan de forma independiente, aunque el sistema ya almacena historiales y resultados en base de datos. Una posible evolución consistiría en permitir conversaciones con contexto acumulado, de forma que el usuario pudiera realizar preguntas encadenadas sobre una misma licitación sin necesidad de repetir información previa.

A nivel de modelos de lenguaje, una posible mejora consistiría en incorporar \textbf{sistemas de inferencia} más eficientes que permitan ejecutar modelos de gran tamaño con menores requisitos \eng{hardware}. Actualmente, el sistema integra Ollama con soporte configurable para CPU y GPU, así como carga dinámica de distintos modelos configurables desde variables de entorno. En este contexto, sería especialmente interesante investigar la integración del proyecto \textit{AirLLM}~\cite{trabfut:airllm}, desarrollado para ejecutar modelos LLM de gran tamaño mediante técnicas de carga dinámica y reducción del consumo de memoria VRAM. Su integración permitiría utilizar modelos más avanzados manteniendo un consumo de recursos reducido, mejorando la escalabilidad y facilitando el despliegue en equipos con recursos limitados.

Otra línea de mejora especialmente relevante sería \textbf{ampliar el sistema de evaluación automática} ya implementado en el proyecto. Actualmente, el sistema incorpora métricas y herramientas de evaluación orientadas al análisis de relevancia y calidad de las respuestas generadas. Como evolución futura, sería especialmente interesante integrar el \eng{framework} \textit{DeepEval}~\cite{trabfut:LLMevaluator}, orientado específicamente a la evaluación y \eng{benchmarking} de aplicaciones basadas en LLM y sistemas RAG. Esto permitiría automatizar pruebas de regresión, comparar modelos y configuraciones de recuperación, detectar degradaciones de calidad y construir un entorno de validación continua para futuras versiones del sistema.

A nivel documental, otra mejora importante consistiría en perfeccionar el procesamiento estructural de los PDF. Durante el desarrollo del proyecto se detectó que la gran heterogeneidad de formatos presentes en los pliegos, dificulta la extracción correcta de secciones y subsecciones. Como línea futura podrían investigarse técnicas multimodales, modelos OCR avanzados o herramientas de análisis estructural capaces de generar documentos Markdown más precisos y estructurados. Esto permitiría enriquecer los metadatos asociados a cada \eng{chunk} y mejorar la recuperación semántica.